Mixture-of-Experts LLMs expose a serving bottleneck that dense-model systems largely avoid: every token chooses a small set of experts, yet the serving engine must keep enough expert weights, caches, and placement state ready before the router's decision is known.
This paper studies whether multi-agent LLM workloads provide earlier signals.
Agent runtimes already know a request's role, phase, tool type, graph position, and prompt-block composition before decoding starts.
\project uses this metadata to predict an agent-conditioned \emph{\EWS} and to drive conservative expert prefetching, scheduling, and placement decisions without changing model outputs.

The key mechanism is \RouteSig, an online distribution store keyed by hierarchical agent metadata.
\RouteSig learns per-layer expert distributions from routed traces, falls back to coarser keys when evidence is sparse, and exposes confidence and miss-cost estimates to the serving policy.
We implement a trace-first prototype with a vLLM-compatible sidecar and a capability gate that disables routed-expert capture on dense models.
On real routed traces from \TinyMixtral, \RouteSig reaches a top-4 expert-label hit rate of 86.1\%, improving over a global-frequency baseline at 77.4\%, and replay simulation estimates a 160.4 ms stall-reduction opportunity with 2.1\% wasted prefetch.
On the requested \Qwen target, a dense Qwen2 model, \project runs end-to-end through real vLLM and correctly takes the no-router fallback path while preserving generated outputs and logging metadata overhead.
Together, these results show both the opportunity for agent-conditioned MoE serving and the systems boundary needed for safe deployment.
